% ============================================================
\section{Introdução e motivação}\label{sec:introducao}
% ============================================================

Transformers representam cada token por um vetor num espaço de alta dimensão. Na entrada do
modelo, esse vetor é uma linha da matriz de \emph{embeddings}: todas as ocorrências de um mesmo
tipo de token começam com \emph{exatamente} o mesmo vetor, não importa o texto em que
apareçam. À medida que atravessam os blocos do Transformer, as representações incorporam
informação do contexto (atenção sobre os tokens anteriores e transformações não lineares) e
se reorganizam no espaço vetorial. No último bloco, a representação de uma ocorrência serve
para prever o \emph{próximo} token, e não mais para identificar o token atual.

A pergunta de fundo deste projeto é \textbf{como} essa reorganização acontece: quanto da
estrutura de similaridade entre representações continua organizada pela identidade lexical e
quanto passa a refletir o contexto, e em que camadas. Estudos anteriores mediram essa
passagem com estatísticas de pares de vetores, como a autossimilaridade das ocorrências de uma
palavra e a anisotropia do espaço~\cite{ethayarajh2019}. Essas medidas resumem distâncias, mas
não dizem \emph{quem é vizinho de quem}, nem se a mudança é só local (troca de alguns vizinhos)
ou se altera a organização em escala maior (\emph{hubs}, comunidades, distâncias).

\subsection{Por que redes complexas}

A proposta é representar cada espaço de representações como uma \textbf{rede}: os vértices são
ocorrências de tokens (ou tipos, na rede lexical de tipos) e cada vértice se liga aos seus $k$
vizinhos mais próximos pelo cosseno. Essa construção tem três vantagens para as perguntas do
projeto:
\begin{itemize}
    \item \textbf{Comparação vértice a vértice.} As redes por ocorrência têm \emph{os mesmos
    vértices} em todas as representações (lexical e camadas 1, 18 e 36). A vizinhança de uma
    ocorrência pode ser comparada diretamente entre camadas (Jaccard), e a fração de vizinhos do
    mesmo tipo (dominância lexical) mede quanto a identidade do token ainda manda.
    \item \textbf{Escalas diferentes.} As ferramentas de redes complexas separam a escala local
    (vizinhança, clusterização) da global (distribuição de graus, \emph{hubs}, distâncias,
    componentes e comunidades). Isso permite distinguir uma reorganização que só troca vizinhos
    de uma que muda a estrutura inteira.
    \item \textbf{Rótulos externos.} As comunidades encontradas podem ser comparadas com
    partições conhecidas (tipo de token, tema, artigo, parágrafo, sentença, próximo token) por
    medidas de concordância, o que transforma ``a rede ficou contextual'' numa afirmação
    mensurável.
\end{itemize}

\subsection{O que o experimento faz}

O experimento usa cerca de $1{,}5\times10^4$ ocorrências de tokens em parágrafos da Wikipédia em
português, escolhidas em oito temas, e o modelo aberto Qwen3-4B-Base~\cite{qwen3} (36 blocos,
dimensão 2560). Para cada ocorrência, extraímos o \emph{embedding} de entrada do seu token e
a saída bruta dos blocos 1, 18 e 36 (e, para a curva camada a camada, dos 36 blocos). Com
esses vetores construímos:
\begin{enumerate}
    \item a \textbf{rede lexical de tipos}, sobre o vocabulário inteiro do modelo (cerca de
    151,6 mil tipos), que descreve a geometria do espaço de entrada;
    \item a \textbf{rede lexical por ocorrência}, em que cada ocorrência recebe o vetor do seu
    tipo, de modo que ocorrências do mesmo tipo são idênticas;
    \item as \textbf{redes contextuais}, com os mesmos vértices representados pelos estados
    internos de cada camada analisada.
\end{enumerate}
As redes são construídas por k-NN com similaridades calculadas em \emph{float64}, com três
tratamentos explícitos de empates (os empates são a regra, e não a exceção, na rede lexical
por ocorrência), nas versões direcionada, por união e mútua. As Seções~\ref{sec:redes}
e~\ref{sec:metodos} descrevem esse método em detalhe.

\subsection{Como ler este relatório}

Este é um \textbf{documento vivo}: ele é atualizado a cada etapa do pipeline e registra o que
foi feito, por que foi feito assim, o que deu errado e o que se encontrou. Ele serve de base
para o relatório final da disciplina, que pode ser uma versão condensada dele. Três convenções
valem para o documento inteiro:
\begin{itemize}
    \item Todo número de resultado citado no texto vem de \cod{tabelas/numeros.tex}, gerado
    pela etapa \etapa{report} a partir dos artefatos da última execução; nada é digitado à
    mão. Enquanto a etapa correspondente não roda, o número aparece como \pendente.
    \item Figuras (\cod{figuras/*.pdf}) e tabelas (\cod{tabelas/*.tex}) também são geradas.
    Quando ainda não existem, um quadro indica o arquivo esperado e a etapa que o produz.
    Números de planejamento (estudo de viabilidade, parâmetros de configuração) são fixos e
    aparecem diretamente no texto, com a fonte indicada.
    \item As Seções~\ref{sec:resultados}--\ref{sec:discussao} ainda são só estrutura: cada
    subseção diz o que vai conter e de qual etapa vêm os dados.
\end{itemize}

\begin{table}[htbp]
\centering
\caption{Estado das seções deste relatório.}\label{tab:estado}
\small
\begin{tabularx}{\linewidth}{@{}l l X@{}}
\toprule
\textbf{Seção} & \textbf{Estado} & \textbf{Depende de} \\
\midrule
1--3. Introdução, perguntas, terminologia & escrita & proposta e plano de execução \\
4. Dados & desenho e viabilidade escritos & números do corpus e da amostra: etapas
\etapa{corpus} e \etapa{sample} \\
5. Modelo e extração & método escrito & diagnósticos: etapa \etapa{extract} \\
6--7. Redes e métodos & método escrito & tempos e estatísticas de empates: etapa \etapa{knn} \\
8. Resultados & estrutura & etapas \etapa{knn}, \etapa{metrics}, \etapa{analyze} e
\etapa{lens} \\
9. Robustez & estrutura & grade de variantes das etapas \etapa{metrics} e \etapa{analyze} \\
10. Discussão & estrutura & resultados \\
11--12. Limitações e problemas & escritas até o planejamento & atualizadas a cada etapa \\
Apêndices & registro de decisões e comandos escritos & tabela de versões: etapa
\etapa{report} \\
\bottomrule
\end{tabularx}
\end{table}
